\documentclass[12pt]{article}
\usepackage[utf8]{inputenc}
\usepackage[T1]{fontenc}
\usepackage[brazilian]{babel}
\usepackage{geometry}
\usepackage{amsmath}
\usepackage{listings}
\usepackage{xcolor}
\usepackage{fancyhdr}
\usepackage{titlesec}
\usepackage{enumitem}

\geometry{margin=1in}
\setlength{\headheight}{15pt}
\setlength{\parindent}{0pt}
\setlength{\parskip}{4pt}

\definecolor{primary}{HTML}{1F4E79}
\definecolor{secondary}{HTML}{0B6E99}
\definecolor{lightbg}{HTML}{F4F8FC}
\definecolor{codebg}{HTML}{F7F7F7}
\definecolor{codeframe}{HTML}{D3D3D3}

\titleformat{\section}
  {\Large\bfseries\color{primary}}
  {}
  {0pt}
  {}

\titleformat{\subsection}
  {\large\bfseries\color{secondary}}
  {}
  {0pt}
  {}

\setlist[itemize]{leftmargin=1.4em,itemsep=2pt,topsep=3pt}
\setlist[enumerate]{leftmargin=1.6em,itemsep=2pt,topsep=3pt}

\pagestyle{fancy}
\fancyhf{}
\lhead{\textbf{Programação em Python}}
% \chead{Aula 01: Algoritmos e Variáveis}
\rhead{Prof. Msc. Nator Junior Carvalho da Costa}
\cfoot{\thepage}
\renewcommand{\headrulewidth}{0.4pt}

\lstset{
    language=Python,
    basicstyle=\ttfamily\small,
    backgroundcolor=\color{codebg},
    frame=single,
    rulecolor=\color{codeframe},
    keywordstyle=\color{blue},
    commentstyle=\color{gray},
    stringstyle=\color{red},
    breaklines=true,
    showstringspaces=false,
    tabsize=4
}

\begin{document}

\begin{center}
{\LARGE \textbf{Lista de Exercícios - Aula 01}}\\[4pt]
{\large Algoritmos e Variáveis}\\[10pt]
\textbf{Disciplina:} Lógica de Programação em Python\\
\textbf{Professor:} Prof. Msc. Nator Junior Carvalho da Costa\\
\textbf{Semestre:} 2026.1
\end{center}

\vspace{8pt}
\colorbox{lightbg}{%
\parbox{\dimexpr\linewidth-2\fboxsep\relax}{%
\textbf{Instruções}\\
Resolva os exercícios abaixo usando Python. Lembre-se:
\begin{itemize}
    \item Um \textbf{algoritmo} é uma sequência de passos para resolver um problema.
    \item Uma \textbf{variável} é um espaço na memória que armazena um valor.
    \item Use nomes descritivos para suas variáveis.
\end{itemize}
}}

\section*{Exercícios}

\subsection*{Exercício 1: Conversão de Temperatura}
Escreva um programa que:
\begin{enumerate}
    \item Crie uma variável \texttt{celsius} com o valor 25
    \item Converta para Fahrenheit usando a fórmula: $F = C \times \frac{9}{5} + 32$
    \item Armazene o resultado em uma variável \texttt{fahrenheit}
    \item Exiba o resultado no formato: \texttt{25°C equivale a 77°F}
\end{enumerate}

\subsection*{Exercício 2: Cálculo de Área}
Escreva um programa que:
\begin{enumerate}
    \item Crie uma variável \texttt{lado} com o valor 5
    \item Calcule a área de um quadrado usando $A = lado^2$
    \item Armazene o resultado em uma variável \texttt{area}
    \item Exiba a mensagem: \texttt{Um quadrado com lado 5 tem área de 25}
\end{enumerate}

\subsection*{Exercício 3: Operações com Variáveis}
Crie um programa que:
\begin{enumerate}
    \item Crie uma variável \texttt{x} com valor 10
    \item Crie uma variável \texttt{y} com valor 3
    \item Calcule e exiba:
    \begin{itemize}
        \item A soma de x e y
        \item A subtração de x por y
        \item A multiplicação de x por y
        \item A divisão de x por y (divisão inteira)
        \item O resto da divisão de x por y
    \end{itemize}
\end{enumerate}

\subsection*{Exercício 4: Custo Total}
Escreva um programa que:
\begin{enumerate}
    \item Crie uma variável \texttt{preco\_unitario} com valor 15.50
    \item Crie uma variável \texttt{quantidade} com valor 4
    \item Calcule o custo total multiplicando preço por quantidade
    \item Armazene em uma variável \texttt{custo\_total}
    \item Exiba a mensagem: \texttt{4 unidades a R\$ 15.50 cada = R\$ 62.00}
\end{enumerate}

\subsection*{Exercício 5: Cálculo de Média}
Escreva um programa que:
\begin{enumerate}
    \item Crie três variáveis: \texttt{nota1 = 8.5}, \texttt{nota2 = 9.0}, \texttt{nota3 = 7.5}
    \item Calcule a média aritmética das três notas
    \item Armazene em uma variável \texttt{media}
    \item Exiba a mensagem: \texttt{A média das notas é 8.33}
\end{enumerate}

\subsection*{Exercício 6: Comprimento da Circunferência}
Escreva um programa que:
\begin{enumerate}
    \item Crie uma variável \texttt{raio} com valor 5
    \item Calcule o comprimento da circunferência usando $C = 2\pi r$
    \item Use \texttt{math.pi} para o valor de $\pi$ (ou use 3.14159)
    \item Armazene em uma variável \texttt{comprimento}
    \item Exiba o resultado com 2 casas decimais
\end{enumerate}

\section*{Desafio Bônus}

\subsection*{Desafio 1: IMC}
Escreva um programa que calcule o Índice de Massa Corporal (IMC):
\begin{enumerate}
    \item Crie variáveis \texttt{peso = 70} (em kg) e \texttt{altura = 1.75} (em metros)
    \item Calcule IMC usando a fórmula: $IMC = \frac{peso}{altura^2}$
    \item Exiba o resultado formatado com 2 casas decimais
\end{enumerate}

\subsection*{Desafio 2: Conversão de Tempo}
Escreva um programa que:
\begin{enumerate}
    \item Crie uma variável \texttt{segundos\_totais = 3665}
    \item Converta para horas, minutos e segundos
    \item Exiba no formato: \texttt{3665 segundos = 1 hora, 1 minuto e 5 segundos}
\end{enumerate}

\end{document}
